% SPDX-License-Identifier: GPL-3.0-or-later
% Copyright (C) 2026 Alessandro Greco
\chapter{意大利和欧盟的法律与伦理问题}\label{cap:25-legale}
\aggiornatoal{2026 年 9 月 23 日}{欧盟条例（EU）2022/2065；意大利第 196/2003 号立法法令（意大利个人数据保护机构公布的合并文本）；Tor Research Safety Board 指南}

\begin{obiettivi}
学完本章后，你将能够：
\begin{itemize}[nosep]
  \item 区分工具本身的合法性和使用工具的行为是否合法；
  \item 根据 Tor 出口中继的运行方式，理解欧盟的\enquote{单纯传输通道}（mere conduit）规则；
  \item 识别意大利中继运营者面临的尚未解决的问题，包括责任、数据保留和设备扣押；
  \item 将 Tor Research Safety Board 的安全原则应用于 Tor 研究。
\end{itemize}
\end{obiettivi}

\begin{attenzione}[本章不构成法律意见]
法律会发生变化，适用结果取决于每个案件的具体事实，而专门涉及 Tor 的司法裁判仍相对有限。本章旨在梳理相关问题，指出已经对照官方文本研究的法律规定，并明确说明不确定性。任何准备作出实际运营决策的人，例如准备运行出口中继的运营者，都应咨询熟悉技术法律的专业律师。
\end{attenzione}

% ---------------------------------------------------------------------------
\section{工具及其用途}\label{sec:legale-strumento}

在意大利以及欧盟其他地区，法律通常并不禁止使用 Tor 或其他保护通信机密性的工具。相反，通信与信函的保密（意大利宪法第~15 条）、表达自由（第~21 条）均受到基本权利保护；《欧盟基本权利宪章》第~7 条和第~8 条同样保护私人生活与个人数据。

人们通过 Tor 进行的活动，仍然受不使用 Tor 时同样适用的法律约束。未经授权访问计算机系统（意大利刑法典第~615-\emph{ter} 条），不会因为流量经过三台中继就变成合法行为。本章的核心是区分工具与其用途；对基础设施提供者也是如此：中继运营者与网络接入服务商类似，提供的是传输能力，而不是选择被传输的内容。

% ---------------------------------------------------------------------------
\section{中继运营者}\label{sec:legale-relay}

\subsection{单纯传输通道（Mere Conduit）}

欧盟《数字服务法》（Digital Services Act，DSA）自 2024 年 2 月 17 日起适用，取代了原先由 2000 年《电子商务指令》建立的中介责任规则。意大利曾通过第 70/2003 号立法法令实施该指令~\cite{dsa2022}。DSA 第~4 条规定了\textbf{单纯传输通道}的责任豁免：对于传输他人提供的信息的服务，如果运营者没有发起传输、没有选择传输对象，也没有选择或修改传输的信息，则在满足适用条件时，无须为该信息承担责任。\cref{fig:conduit} 从 cell 协议层面，将这些条件与 Tor 出口中继的行为进行比较。

\begin{figure}[htbp]
\centering
% SPDX-License-Identifier: GPL-3.0-or-later
% Copyright (C) 2026 Alessandro Greco
%% EPUB 无障碍替代文本。
\altfig{本图仅比较 DSA 第 4 条与出口中继的技术行为，不自动确定法律免责。客户端通过 RELAY\_BEGIN 指定目的地，出口建立 TCP 并执行策略，将 RELAY\_DATA 中的应用字节写入 TCP 流，而非原样转发 Tor cell。}
% 单纯传输通道条件与出口中继行为对比
\begin{tikzpicture}[font=\sffamily\scriptsize,node distance=0.35cm,
  cond/.style={draw,rounded corners=2pt,text width=4.2cm,align=left,minimum height=0.8cm},
  tor/.style={draw=torviola,fill=torviola!8,rounded corners=2pt,text width=5.4cm,align=left,minimum height=0.8cm}]
  \node[cond] (a) {（a）不发起传输};
  \node[cond,below=of a] (b) {（b）不选择接收方};
  \node[cond,below=of b] (c) {（c）不选择或修改信息};
  \node[tor,right=1.2cm of a] (ta) {客户端发送 \texttt{RELAY\_BEGIN}；出口打开 TCP};
  \node[tor,right=1.2cm of b] (tb) {目的地由客户端指定；策略可拒绝};
  \node[tor,right=1.2cm of c] (tc) {\texttt{RELAY\_DATA} 字节写入 TCP 流};
  \foreach \x in {a,b,c} \draw[<->,torgrigio] (\x) -- (t\x);
  \node[above=4pt of a,font=\sffamily\small] {欧盟条例（EU）2022/2065，第~4 条};
  \node[above=4pt of ta,font=\sffamily\small] {Tor 出口中继};
  \node[font=\sffamily\tiny,text=torgrigio,align=center,below=0.5cm of c,xshift=3.3cm] {免责须逐案判断};
\end{tikzpicture}

\caption{示意比较欧盟条例（EU）2022/2065 第~4 条的\emph{单纯传输}条件与 Tor 出口的技术行为，不构成自动免责结论。客户端通过 \texttt{RELAY\_BEGIN} 指定目的地；出口建立 TCP 连接，并可按照出口策略拒绝请求。\texttt{RELAY\_DATA} 中的应用数据字节写入 TCP 流，而不是将 Tor cell 原样转发。是否免责还取决于全部法律条件与具体情况。}
\label{fig:conduit}
\end{figure}

还有两项规定值得关注。第~8 条排除了对所传输信息进行\emph{一般性}监控，或主动寻找违法行为证据的义务：运营者通常没有无差别检查流量的法律义务；中间中继也无法读取被加密的应用层明文内容。另一方面，第~4 条第 3 款保留了司法或行政主管机关命令服务提供者预防或终止侵权行为的权力。

\paragraph{尚未解决的问题。}该条例适用于\enquote{信息社会服务}。按照欧盟法律定义，此类服务通常是\emph{有偿提供}的。完全免费、由志愿者运营的中继是否属于这一范畴？一种观点主张应作较宽泛的解释，而且责任豁免反映了法院可能考虑的原则：仅负责传输信息的人，不应为信息内容承担责任。但是，本书并未找到能够直接确立意大利志愿者 Tor 中继法律地位的适用判决。这种不确定性也有助于理解，为什么有些出口中继由本来就提供网络服务的协会或组织运营。

\subsection{民事责任与刑事调查}

豁免规则涉及对所传输信息承担的责任，却不能阻止刑事调查从出口 IP 地址开始。受害者看到的通过 Tor 实施的违法行为，表面上来自出口中继。调查人员在初期可能搜查或扣押运营者的计算机。Electronic Frontier Foundation 为中继运营者编写的法律常见问题解答虽以美国法为背景，但也总结了实务经验：该组织表示不知道有人仅因运行中继而被判有罪，却记载了由于误解 Tor 而开展的调查与扣押~\cite{efflegalfaq}。以下措施可以降低实际风险：
\begin{itemize}
  \item 在数据中心的专用服务器上运行出口，不要使用家用计算机或个人工作电脑。
  \item 清楚表明该 IP 地址属于出口中继，例如设置说明性的反向 DNS，或使用 \opt{DirPortFrontPage} 在 tor 的 DirPort 上提供介绍页面。
  \item 事先通知网络接入服务商，并约定如何处理滥用投诉。
  \item 回应投诉时解释 Tor 的工作方式，并保存相关通信记录。
\end{itemize}

\subsection{数据保留}

意大利《个人数据保护法典》第~132 条要求受其约束的\emph{服务提供者}为刑事调查和起诉保留电子通信流量数据十二个月，而电话流量数据的保留期为二十四个月~\cite{codiceprivacy}。第~5-\emph{bis} 款保留了特殊规定：2017 年 11 月 20 日第~167 号法律第~24 条，对于《刑事诉讼法典》第~51 条第 3-\emph{quater} 款和第~407 条第 2 款（a）项列明的罪行——包括恐怖主义及其他严重犯罪——将调查所需的保留期限延长为\textbf{七十二个月}~\cite{legge167}。《个人数据保护法典》的相关部分适用于\emph{通过公共通信网络向公众提供的}电子通信服务（第~121 条）。

志愿者中继运营者是否属于此类服务提供者，仍然是一个未解决的问题，并且可能取决于服务的具体提供方式。运营者，特别是原本就提供通信服务的组织，在运行出口前应请律师评估适用情况。从技术上说，在推荐的日志设置下，tor 通常不会保存用户目的地的历史记录。出口中继能看到它所转发连接的目的地，却不知道原始客户端的 IP 地址。实际能够获得的数据取决于中继的角色和配置，因此不能假定出口中继可以识别 Tor 的具体用户。

\subsection{日志属于个人数据}

IP 地址属于个人数据。《通用数据保护条例》（GDPR）要求仅处理必要的数据，并且仅在必要的期限内保留（第~5 条，数据最小化原则）。对于中继而言，符合安全目标的选择是仅以 \texttt{notice} 级别记录日志，保持 \opt{SafeLogging} 启用，并避免抓包（\cref{sec:traffico-log}）。

% ---------------------------------------------------------------------------
\section{研究人员}\label{sec:legale-ricerca}

Tor 研究历史上曾出现一些存在伦理争议的实验，部分案例将在 \cref{cap:27-casi-rete} 中讨论。有些研究人员运行中继，监控或操纵真实用户的流量，以展示某种攻击。这些实验除了产生伦理问题，还可能触及意大利刑法，包括关于非法截获计算机通信的第~617-\emph{quater} 条，以及个人数据保护方面的规定。

\textbf{Tor Research Safety Board} 是一个为实验安全提供建议的研究人员组织，并制定了相关指南~\cite{trsb}：
\begin{itemize}
  \item 尽可能使用测试网络，只攻击自己产生的流量。
  \item 仅收集真正需要的数据，并且只收集适合安全公开的数据。
  \item 保护收集到的数据以及实验工具。
  \item 通过舍入或加入噪声降低数据精度。
  \item 考虑第三方数据集与自己收集的数据组合后可能带来的信息泄露。
  \item 评估用户是否有理由认为相关数据应保持私密。
  \item 确保预期收益足以证明风险是合理的。
\end{itemize}
安全委员会要求研究人员说明希望了解什么、准备如何实施、实验会引入哪些风险，以及为什么预期收益能够抵消风险。本书遵循同样的原则：需要产生实际 Tor 流量的实验使用私有测试网络（\cref{sec:rp-perche}）；而 \cref{cap:24-metrics} 中的公共文档分析则离线进行，使用公开的汇总数据。

% ---------------------------------------------------------------------------
\section*{本章小结}
\begin{itemize}
  \item 使用 Tor 本身是合法的；通过 Tor 开展的活动仍受一般法律约束。
  \item DSA 第~4 条在满足规定条件时，为单纯传输他人信息的服务提供责任豁免。出口中继的运行方式与这些条件有关，但志愿者中继能否适用该条例，仍存在法律不确定性。
  \item 责任豁免不能阻止调查或扣押。使用专用服务器、明确标识出口身份，以及与网络接入服务商协调，可以减少实际问题。
  \item 数据保留要求适用于法律涵盖的、向公众提供的通信服务提供者；这些要求是否以及如何适用于志愿者中继，需要逐案进行法律分析。
  \item Tor 研究应尽可能避免涉及真实用户的实验：使用测试网络、最小化数据、只发布安全的汇总数据，并权衡收益与风险。
\end{itemize}

% ---------------------------------------------------------------------------
\section*{练习}

\begin{esercizio}[1]\label{es:25-condizioni}
针对 DSA 第~4 条的三项条件，分别描述一种出口中继可能采取的、会违反该条件的行为变化。
\end{esercizio}

\begin{esercizio}[1]\label{es:25-strumento}
通过一个实例解释工具合法性与使用行为合法性的区别，并分别将其应用于 Tor 用户和中继运营者。
\end{esercizio}

\begin{esercizio}[2]\label{es:25-segnalazione}
假设某网站投诉你的出口 IP 地址发起了多次访问尝试，网络接入服务商因此向你发送滥用报告。请草拟一份回应，说明你能够解释哪些情况，以及可以采取哪些实际措施。
\end{esercizio}

\begin{esercizio}[2]\label{es:25-esperimento}
某研究人员提议记录出口中继观察到的域名，以估计有多少 Tor 用户访问意大利新闻网站。请根据 Tor Research Safety Board 的指南评估这一计划，并提出更安全的替代方案。
\end{esercizio}

\begin{esercizio}[3]\label{es:25-ricerca}
查找意大利及其他欧洲国家法院涉及 Tor 中继运营者的司法裁判。分别总结争议的法律问题与结果，并引用一手来源。
\end{esercizio}

% ---------------------------------------------------------------------------
\section*{延伸阅读}
\begin{itemize}
  \item 欧盟条例（EU）2022/2065，尤其是第~4、8 和~9 条~\cite{dsa2022}。
  \item 意大利第 196/2003 号立法法令第~121 条及第~132 条，以意大利个人数据保护机构发布的合并文本为准~\cite{codiceprivacy}。
  \item Electronic Frontier Foundation，\titolo{Legal FAQ for Tor Relay Operators}~\cite{efflegalfaq}。
  \item Tor Research Safety Board 网站~\cite{trsb}。
\end{itemize}
